\documentclass[10pt,a4paper]{amsart}
\usepackage[margin=19mm]{geometry}
\usepackage{amsmath,amssymb,mathtools}
\usepackage[scr=rsfs,cal=euler]{mathalfa}
\usepackage{xfrac}
\usepackage[unicode,hidelinks,bookmarks=false]{hyperref}
\newcommand{\ie}{i.e.\ }
\newcommand{\cf}{cf.\ }
\newcommand{\resp}{respectively\ }
\newcommand{\Subsection}{\subsection}
\providecommand{\nobreakdash}{\nobreak\hbox{-}}
\setlength{\emergencystretch}{2em}
\multlinegap0pt
\usepackage{amssymb}
\usepackage{mathtools}
\usepackage[shortlabels]{enumitem}
\usepackage{xcolor}
\newtheorem{claim}{Claim}
\newtheorem{lemma}{Lemma}
\title{Two conjectures on graphs and their edge-path matrices}
\author{Metrose Metsidik, Xian'an Jin}
\date{Source: arXiv:2608.15992v1 (CC BY 4.0)}
\begin{document}
\maketitle

\noindent\emph{Source and licence.} Two conjectures on graphs and their edge-path matrices. Version arXiv:2608.15992v1 (CC BY 4.0), distributed by arXiv under the Creative Commons Attribution 4.0 International licence, http://creativecommons.org/licenses/by/4.0/, as recorded in the arXiv licence metadata for this exact version and shown on the abstract page https://arxiv.org/abs/2608.15992v1. Full licence notice: \texttt{licenses/arxiv-2608.15992-CC-BY-4.0.txt}.\par\smallskip

\noindent\textit{Editorial scope.} Counted unit: the article's lemma ML-1002 (source0.tex lines 356--358). The article's complete original proof of it is delivered with it (source0.tex lines 359--569, one \texttt{proof} environment of 211 lines). No mathematics is written, repaired or re-derived: the statement and the proof are the article's own lines, in the article's own notation, and no retained line is substituted or re-flowed.  Retained uncounted as the source-local context the proof needs: lines 232--241.  Nothing was omitted from any retained span, so the package carries no omission record.  No retained line contains a citation, so the wrapper carries no bibliography and no bibliography entry.  No retained line contains a figure environment, an image-inclusion command or drawing code, so no image is delivered and the package declares no asset dependency.  Editorial formatting only, in addition: an AMS \texttt{amsart} preamble that restates the article's own declarations and macros used by the retained lines, namely the environments \texttt{claim}, \texttt{lemma} on separate counters, together with the standard packages those lines need.  The retained environments print in this document as Claim 1, Lemma 1 in order of appearance, whereas the article prints the same items with the numbering of its own full document.  The complete native source of the article as distributed in the arXiv e-print for this exact version is bundled unchanged as \texttt{sources/207801/source0.tex}.
\begin{claim}
\label{cl:stronger}
\[
|\partial X|\geq
\begin{cases}
2\sum\limits_{s\in S^+}\min\{|X|,s\}, & \text{if } n/2\notin S,\\[4pt]
2\sum\limits_{s\in S^+\setminus\{n/2\}}\min\{|X|,s\}+\min\{|X|,n/2\}, & \text{if } n/2\in S.
\end{cases}
\]
\end{claim}

\begin{lemma}\label{ML-1002}
Let \(G=\operatorname{Cir}(n,S)\), where \(q=|S|\), \(n\neq 6\), \(l\in S\), and, if \(n\) is even, \(r\in S\). Suppose \(H\subseteq G\) satisfies \(|E(H)|<n\), \(\Delta(H)\le 2\), and \(H\) is not \(2\)-regular. Then, for every \(t\in\{q-1,q\}\), any two vertices \(u,v\) whose degrees in \(G-E(H)\) are at least \(t\) are joined by \(t\) edge-disjoint \(u\)--\(v\) paths.
\end{lemma}

\begin{proof}
The proof of this lemma follows a similar idea to that of the previous lemma, but the computation requires more care. 

Let  
\[
U=\{u:\deg_{G-E(H)}(u)\ge t\}.
\]
Suppose that for some \(u,v\in U\), there are fewer than \(t-1\) edge-disjoint \(u\)-\(v\) paths in \(G-E(H)\). By Menger's theorem, there exists an edge cut \(C\subseteq E(G-E(H))\) separating \(u\) and \(v\) with \(|C|\le t-2\). Choose such a cut of minimum cardinality, and let \(X\) be the vertex set of the component of \(G-C\) containing \(u\). Set $\overline{X}=V(G)\setminus X$. Then
\[
C=\partial_{G-E(H)}X=\partial_{G-E(H)}\overline{X}.
\]
Since \(\deg_{G-E(H)}(u)\ge t\) and \(\deg_{G-E(H)}(v)\ge t\), we have \(|X|\ge 2\) and \(|\overline{X}|\ge 2\). If \(|X|=2\), then
\[
|\partial_{G-E(H)}(X)|\ge
\begin{cases}
q-1+q-3\ge q, & \text{if } t=q,\\[4pt]
q-2+q-3\ge q-1, & \text{if } t=q-1,
\end{cases}
\]
a contradiction. If \(|X|=3\) and \(q>4\), then
\[
|\partial_{G-E(H)}(X)|\ge
\begin{cases}
q-2+2(q-4)\ge q, & \text{if } t=q,\\[4pt]
q-3+2(q-4)\ge q-1, & \text{if } t=q-1,
\end{cases}
\]
again a contradiction. If \(|X|=3\) and \(q=4\), then since \(n>6\), we have \(l\ge 3\), and therefore, by Claim~\ref{cl:stronger},
\[
|\partial_{G-E(H)}(X)|\ge
\begin{cases}
8-4\ge q, & \text{if } t=q,\\[4pt]
8-5\ge q-1, & \text{if } t=q-1,
\end{cases}
\]
once more a contradiction. Hence \(|X|\ge 4\) and \(|\overline{X}|\ge 4\), so \(n=|X|+|\overline{X}|\ge 8\). 

Now let \(B=\partial_G X\cap H\) be the set of edges of \(H\) crossing the cut. Then
\[
\partial_G X=C\cup B,\qquad C\cap B=\varnothing,
\]
so \(|\partial_G X|=|C|+|B|\). Define \(H_X=H\cap E(G[X])\) and \(U_X=U\cap X\). Clearly \(|U_X|\ge 1\) because \(u\in U_X\).

Each edge of \(H_X\) covers at least one vertex of \(X\), and the endpoints in \(X\) of the edges of \(B\) account for at least \(\lceil |B|/2\rceil\) vertices of \(X\). Therefore,
\[
|X|\ge
\begin{cases}
|U_X|+|H_X|+\lceil |B|/2\rceil, & \text{if } \deg_{G-E(H)}(u)=q,\\[8pt]
|U_X|+|H_X|+\lceil (|B|-1)/2\rceil, & \text{if } \deg_{G-E(H)}(u)=q-1.
\end{cases}
\]
Since \(|U_X|\ge 1\), we obtain
\[
|B|\le
\begin{cases}
2|X|-3, & \text{if } |B| \text{ is odd and } \deg_{G-E(H)}(u)=q,\\[8pt]
2|X|-2, & \text{otherwise}.
\end{cases}
\]
Consequently,
\[
|C|=|\partial_G X|-|B|\ge
\begin{cases}
|\partial_G X|-2|X|+3, & \text{if } |B| \text{ is odd and } \deg_{G-E(H)}(u)=q,\\[8pt]
|\partial_G X|-2|X|+2, & \text{otherwise}.
\end{cases}
\tag{2}
\]

As before, we discuss the cases according to the parity of \(n\). Using Claim~\ref{cl:stronger} and Inequality~(2), we derive a lower bound for \(|C|\).

{\bf Cace A}. $n$ is odd.
\begin{eqnarray*}
|C| &\geq& |\partial_G X|-2|X|+2 \\
&\geq &2+2|X|+2\sum\limits_{s\in S^+\setminus\{s_1,l\}}\min\{|X|,s\}-2|X|+2\\
&\geq &4+4(q/2-2)\\
&\geq &q
\end{eqnarray*}

{\bf Cace B}. $n$ is even.

{\bf Cace B.1}  $|X|\le l$.
\begin{align}
|C| &\geq 
\begin{cases}
\displaystyle 6 + 2\sum_{s\in S^+\setminus\{r,l\}} \min\{|X|,s\}, & \text{if } n/2\notin S \\[4pt]
\displaystyle 6 + 2\sum_{s\in S^+\setminus\{r,l,n/2\}} \min\{|X|,s\}+ |X|, & \text{if } n/2\in S
\end{cases} \\
&\geq 
\begin{cases}
\displaystyle 6 + 2\times
\begin{cases}
0, & \text{if } q=4 \\
1, & \text{if } q=6 \\
3, & \text{if } q=8 \\
2(q/2-2), & \text{if } q\ge 10
\end{cases}, & \text{if } n/2\notin S \\[4pt]
\displaystyle 6+ |X| + 2\times
\begin{cases}
0, & \text{if } q=5 \\
1, & \text{if } q=7 \\
3, & \text{if } q=9 \\
2\big((q-1)/2-2\big), & \text{if } q\ge 11
\end{cases}, & \text{if } n/2\in S
\end{cases} \\
&\geq 
\begin{cases}
\displaystyle q+2, & \text{if } n/2\notin S \\[4pt]
\displaystyle q+5, & \text{if } n/2\in S
\end{cases}
\end{align}
where $r\in\{2,3\}$.

{\bf Cace B.2}  $|X|- l=1$.
\begin{align*}
|C| &\geq 
\begin{cases}
\displaystyle 4 + 2l+2\sum_{s\in S^+\setminus\{2,l\}} \min\{|X|,s\}-2|X|+2, & \text{if } n/2\notin S \\[4pt]
\displaystyle 4 +2l+ 2\sum_{s\in S^+\setminus\{2,l,n/2\}} \min\{|X|,s\}+ l-2|X|+2, & \text{if } n/2\in S
\end{cases} \\
&\geq 
\begin{cases}
\displaystyle 4 + 2\times
\begin{cases}
0, & \text{if } q=4 \\
1, & \text{if } q=6 \\
3, & \text{if } q=8 \\
2(q/2-2), & \text{if } q\ge 10
\end{cases}, & \text{if } n/2\notin S \\[4pt]
\displaystyle 4+ l + 2\times
\begin{cases}
0, & \text{if } q=5 \\
1, & \text{if } q=7 \\
3, & \text{if } q=9 \\
2\big((q-1)/2-2\big), & \text{if } q\ge 11
\end{cases}, & \text{if } n/2\in S
\end{cases} \\
&\geq 
\begin{cases}
\displaystyle q, & \text{if } n/2\notin S \\[4pt]
\displaystyle q+2, & \text{if } n/2\in S
\end{cases}
\end{align*}

{\bf Cace B.3}  $|X|- l=2$.

Since $n \equiv 2 \pmod{4}$, we have $n \ge 10$.

{\bf Cace B.3.1}  $n>10$.

Then $l \ge 5$.
\begin{align*}
|C| &\geq 
\begin{cases}
\displaystyle 6 + 2l+2\sum_{s\in S^+\setminus\{3,l\}} \min\{|X|,s\}-2|X|+2, & \text{if } n/2\notin S \\[4pt]
\displaystyle 6 +2l+ 2\sum_{s\in S^+\setminus\{3,l,n/2\}} \min\{|X|,s\}+ l-2|X|+2, & \text{if } n/2\in S
\end{cases} \\
&\geq 
\begin{cases}
\displaystyle 4 + 2\times
\begin{cases}
0, & \text{if } q=4 \\
1, & \text{if } q=6 \\
3, & \text{if } q=8 \\
2(q/2-2), & \text{if } q\ge 10
\end{cases}, & \text{if } n/2\notin S \\[4pt]
\displaystyle 4+ l + 2\times
\begin{cases}
0, & \text{if } q=5 \\
1, & \text{if } q=7 \\
3, & \text{if } q=9 \\
2\big((q-1)/2-2\big), & \text{if } q\ge 11
\end{cases}, & \text{if } n/2\in S
\end{cases} \\
&\geq 
\begin{cases}
\displaystyle q, & \text{if } n/2\notin S \\[4pt]
\displaystyle q+2, & \text{if } n/2\in S
\end{cases}
\end{align*}

{\bf Cace B.3.2}  $n=10$.
\begin{align*}
|C| &\geq 
\begin{cases}
\displaystyle 8 + 6+2\sum_{s\in S^+\setminus\{3,4\}} s-10+2, & \text{if } 5\notin S \\[4pt]
\displaystyle 8 +6+ 2\sum_{s\in S^+\setminus\{3,4,5\}} s+ 3-10+2, & \text{if } 5\in S
\end{cases} \\
&\geq 
\begin{cases}
\displaystyle 6 + 2\times
\begin{cases}
0, & \text{if } q=4 \\
1, & \text{if } q=6 \\
3, & \text{if } q=8 
\end{cases}, & \text{if } 5\notin S \\[4pt]
\displaystyle 6+ 3 + 2\times
\begin{cases}
0, & \text{if } q=5 \\
1, & \text{if } q=7 \\
3, & \text{if } q=9 
\end{cases}, & \text{if } 5\in S
\end{cases} \\
&\geq 
\begin{cases}
\displaystyle q+2, & \text{if } 5\notin S \\[4pt]
\displaystyle q+4, & \text{if } 5\in S
\end{cases}
\end{align*}
Thus, in every case we obtain \(|C| \ge t\), which contradicts the assumption \(|C| \le t-1\).
\end{proof}

\end{document}
